% FILE: main.tex

\section*{Цель работы}
Изучение принципов построения и функционирования синхронных триггеров (D-типа и JK-типа) на языке описания аппаратуры Verilog. Анализ их аппаратной реализации (RTL) и временных диаграмм работы, включая исследование влияния сигналов управления, таких как синхронный сброс.

\subsection*{ \fbox{ Задание 1 и 1а. } }
\begin{quote}
    \textit{Собрать схему синхронного D-триггера в текстовом редакторе. Изучить схему, реализованную в RTL-Viewer. Построить временные диаграммы, иллюстрирующие работу устройства. Период тактового сигнала задать $20$ нс, время моделирования --- $500$ нс. }

    \textit{Ввести в D-триггер сигнал сброса. Проделать всё то же самое, что и для простого D-триггера. Синхронным или асинхронным является сигнал сброса? Объяснить, почему.}
\end{quote}
 
\subsubsection*{Реализация на языке Verilog}
Коды на языке Verilog для D-триггера и D-триггера с сигналом сброса представлены в листингах \cref{lst:d_trigger} и \cref{lst:d_trigger_reset} соответственно.

\begin{listing}[H]
    \caption{Verilog-описание D-триггера (d.v)}
    \label{lst:d_trigger}
    \inputminted[]{verilog}{scripts/d.v}
\end{listing}

\begin{listing}[H]
    \caption{Verilog-описание D-триггера с сигналом сброса (d+r.v)}
    \label{lst:d_trigger_reset}
    \inputminted[]{verilog}{scripts/d+r.v}
\end{listing}

\subsubsection*{RTL-схемы}
Аппаратные реализации, синтезированные средой разработки, показаны на \cref{fig:d_rtl}.

\begin{figure}[H]
    \centering
    \begin{subfigure}[b]{0.48\textwidth}
        \centering
        \input{figs/d_RTL.pdf}
        \caption{D-триггер}
        \label{fig:d_rtl_simple}
    \end{subfigure}
    \hfill
    \begin{subfigure}[b]{0.48\textwidth}
        \centering
        \input{figs/dr_RTL.pdf}
        \caption{D-триггер с сигналом сброса}
        \label{fig:d_rtl-reset}
    \end{subfigure}
    \caption{RTL-схемы D-триггеров}
    \label{fig:d_rtl}
\end{figure}

\subsubsection*{Временные диаграммы}
Временные диаграммы, иллюстрирующие работу триггеров, приведены на \cref{fig:d_wave}.

\begin{figure}[H]
    \centering
    \begin{subfigure}[b]{0.48\textwidth}
        \centering
        \includegraphics[width=\textwidth]{figs/d_wave_t.png}
        \caption{D-триггер}
        \label{fig:d_wave-simple}
    \end{subfigure}
    \hfill
    \begin{subfigure}[b]{0.48\textwidth}
        \centering
        \includegraphics[width=\textwidth]{figs/dr_wave_t.png}
        \caption{D-триггер с сигналом сброса}
        \label{fig:d_wave_reset}
    \end{subfigure}
    \caption{Временные диаграммы работы D-триггеров}
    \label{fig:d_wave}
\end{figure}

\subsubsection*{Анализ сигнала сброса}
Сигнал сброса \mintinline{verilog}§reset§ является \textit{синхронным}.

\textit{Обоснование:}
Анализ Verilog-кода в \cref{lst:d_trigger_reset} показывает, что проверка условия \mintinline{verilog}§if (reset)§ находится внутри событийного блока \mintinline{verilog}§always @(posedge clock)§. Это означает, что состояние входа \mintinline{verilog}§reset§ опрашивается и влияет на выход \mintinline{verilog}§q§ не в произвольный момент времени, а строго по переднему фронту тактового сигнала \mintinline{verilog}§clock§.

Данное поведение должно быть отражено на временной диаграмме (\cref{fig:d_wave_reset}): изменение состояния выхода \mintinline{verilog}§q§ в ответ на сигнал \mintinline{verilog}§reset§ происходит не мгновенно при его активации, а только по приходу следующего активного тактового фронта. Асинхронный сброс потребовал бы включения сигнала \mintinline{verilog}§reset§ в список чувствительности, например: \mintinline{verilog}§always @(posedge clock or posedge reset)§.

\subsection*{ \fbox{ Задание 2. } }
\begin{quote}
    \textit{Собрать схему синхронного J-K-триггера в текстовом редакторе. Изучить схему, реализованную в RTL-Viewer. Построить временные диаграммы, иллюстрирующие работу устройства. Период тактового сигнала задать 40 нс, время моделирования – 900 нс.}
\end{quote}

\subsubsection*{Реализация и RTL-схема}
Verilog-описание JK-триггера представлено в \cref{lst:jk_trigger}. Соответствующая ему RTL-схема приведена на \cref{fig:jk_rtl}.

\begin{listing}[H]
    \caption{Verilog-описание JK-триггера (jk1.v)}
    \label{lst:jk_trigger}
    \inputminted[]{verilog}{scripts/jk1.v}
\end{listing}

\begin{figure}[H]
    \centering
    \input{figs/jk_RTL.pdf}
    \caption{RTL-схема JK-триггера}
    \label{fig:jk_rtl}
\end{figure}

\subsubsection*{Временная диаграмма и режимы работы}
Временная диаграмма, представленная на \cref{fig:jk_wave}, демонстрирует все четыре режима работы JK-триггера:
\begin{enumerate}
    \item \textit{Режим хранения.} При $J=0, K=0$ состояние триггера не изменяется с приходом тактового импульса.
    \item \textit{Режим сброса (Reset).} При $J=0, K=1$ триггер переходит в состояние $q=0$ по активному фронту тактового сигнала.
    \item \textit{Режим установки (Set).} При $J=1, K=0$ триггер переходит в состояние $q=1$.
    \item \textit{Инверсный (счетный) режим.} При $J=1, K=1$ триггер инвертирует свое состояние при каждом тактовом импульсе.
\end{enumerate}

\begin{figure}[H]
    \centering
    \includegraphics[width=\textwidth]{figs/jk_wave_t.png}
    \caption{Временная диаграмма работы JK-триггера}
    \label{fig:jk_wave}
\end{figure}

\section*{Выводы}
В ходе выполнения работы были реализованы и проанализированы HDL-модели синхронных D- и JK-триггеров. Были построены их аппаратные схемы (RTL) и временные диаграммы, иллюстрирующие их функционирование. Проведено исследование влияния сигнала синхронного сброса на работу D-триггера. Анализ временных диаграмм подтвердил теоретические принципы работы данных цифровых устройств.
